\chapter{Серия \texttt{6x}: от topology correction к гибриду local + global}

\section{Зачем вообще понадобилась серия \texttt{6x}}

После \texttt{exp4} стало ясно, что семантический gate сам по себе не вытягивает
специфический ответ из-под широкого хаба. Но это еще не говорило, \emph{какой
именно} механизм отсутствует. Серия \texttt{6x} строится как последовательное
расщепление вопроса:

\begin{enumerate}[leftmargin=2em]
\item действительно ли дело в топологии сходимости путей;
\item если да, можно ли исправить это локальной поправкой на удержание;
\item если локальная поправка помогает, какой именно класс знаменателей
осмыслен;
\item если локального удержания мало, нужен ли еще и глобальный structural
prior.
\end{enumerate}

Важно, что здесь впервые появляется уже почти ``физическая'' интерпретация
узла: не всякий inflow одинаково превращается в salience, потому что разные узлы
по-разному удерживают и рассеивают приходящий поток.

\section{\texttt{exp6}: topology-aware retrieval correction}

\subsection{Цель опыта}

\texttt{exp6} должен был сделать то, что \texttt{exp4} еще не мог: отделить
собственно semantic relevance от эффекта многократной сходимости на хабе.

Для этого опыт делился на две части:

\begin{enumerate}[leftmargin=2em]
\item topology surgery: меняется только структура путей;
\item correction ablation: исходный граф сохраняется, меняется только формула
поправки.
\end{enumerate}

\subsection{Часть A: хирургия топологии}

В первой части проверялась почти причинная гипотеза:

\begin{quote}
если проблема действительно в количестве и силе сходящихся маршрутов, то
достаточно удалить пару входов в \texttt{beverage} или добавить пару входов в
\texttt{caffeine}, и ranking должен перевернуться даже без смены embeddings и
query.
\end{quote}

Именно это и произошло.

\begin{center}
\begin{tabular}{lccc}
\toprule
Вариант графа & \texttt{caffeine} & \texttt{beverage} & Победа \\
\midrule
Исходный & 27.20 & 129.46 & нет \\
Срезать часть входов в \texttt{beverage} & 37.57 & 27.38 & да \\
Добавить входы в \texttt{caffeine} & 95.10 & 77.17 & да \\
Сделать оба изменения & 141.30 & 19.26 & да \\
\bottomrule
\end{tabular}
\end{center}

Это один из самых чистых результатов всей серии: при неизменном query и
неизменном semantic gate ranking переворачивается одной лишь правкой topology.

\subsection{Часть B: пост-аккумуляционные поправки}

После этого можно было уже предметно тестировать корректирующие формулы. На
исходном графе сравнивались:

\begin{itemize}[leftmargin=2em]
\item отсутствие поправки;
\item деление на $\deg(B)^\gamma$;
\item деление на $\fanin_t(B)^\gamma$.
\end{itemize}

Ключевым было то, что штраф накладывался \emph{после слияния вкладов} в узле
назначения. Это принципиально: если наказывать поток локально у источника, то
узел, у которого есть единственный путь наружу, все равно отправит весь свой
заряд в этот путь.

\subsection{Результат}

Единственной формулой, действительно перевернувшей ranking, оказалось деление
на полную степень узла:

\[
\incoming(B)\mapsto \frac{\incoming(B)}{\deg(B)}.
\]

Это был грубый, но убедительный proof of concept.

\subsection{Что эксперимент реально доказал}

\texttt{exp6} не доказал, что $1/\deg(B)$ есть закон памяти. Он доказал более
узкое, но очень важное утверждение:

\begin{quote}
если узел побеждает за счет чрезмерной сходимости путей, то поправка должна
действовать на merged accumulation, а не только на локальный выбор ребра у
источника.
\end{quote}

\section{\texttt{exp6b}: richer golden world и первое сравнение семейств}

\subsection{Почему одного кофейного fixture уже было мало}

Успех на одной конфигурации \texttt{caffeine} против \texttt{beverage} еще не
говорит, что найдена хорошая retrieval law. Можно легко подобрать формулу,
которая работает только на одном сценарии.

Поэтому \texttt{exp6b} строит богатый, но все еще читаемый золотой мир на 33
узлах. Его смысл не в ``реализме мира'', а в разнообразии структурных ситуаций.

\subsection{Четыре кластера графа}

В графе были выделены:

\begin{itemize}[leftmargin=2em]
\item домашний кластер (\texttt{apartment}, \texttt{wallet}, \texttt{keys},
\texttt{cats}, \texttt{feeding alarm}, \texttt{elevator turnaround});
\item travel cluster (\texttt{trip}, \texttt{airport}, \texttt{passport},
\texttt{boarding pass}, \texttt{hotel});
\item work cluster (\texttt{office}, \texttt{meeting}, \texttt{laptop},
\texttt{charger}, \texttt{dead laptop});
\item coffee/health cluster (\texttt{coffee}, \texttt{caffeine},
\texttt{beverage}, \texttt{energy}, \texttt{water}).
\end{itemize}

Принципиально важно, что запросы включали не только anti-hub cases, но и
legitimate-hub cases. Иначе любую агрессивную нормировку можно было бы ошибочно
объявить хорошей.

\subsection{Формулы}

Появляются семейства

\[
R_0(B)=1,
\]
\[
R_1(B)=1+\beta\, \degout(B)^\gamma,
\]
\[
R_2(B)=1+\beta\, \sout(B)^\gamma,
\]
\[
R_3(B)=\alpha+\beta\, \sout(B)^\gamma,
\]
\[
R_4(B)=\capacity(B)+\beta\, \sout(B)^\gamma,
\]
\[
R_5(B)=\alpha+\capacity(B)+\beta\, \sout(B)^\gamma.
\]

Здесь впервые появляется capacity $C(B)$ как локальная ширина представления.

\subsection{Что можно было ожидать}

Если ``утечка'' исчерпывает всю проблему, должен был победить leak-only класс
$R_2$ или $R_3$. Если же разные узлы отличаются не только по суммарной силе
исходящих каналов, но и по собственной концептуальной ширине, тогда capacity
должна дать добавочное преимущество.

\subsection{Результат}

В \texttt{exp6b} лучшим in-sample и leave-one-out победителем стало семейство

\[
R_4(B)=\capacity(B)+0.1\, \sout(B)^{0.5}.
\]

Оно решило все 7 запросов и дало наибольший общий margin. Базовый raw case
проваливал только 6/7, а leak-only формулы не смогли стать убедительными
победителями.

\subsection{Интерпретация}

Это был первый сильный аргумент в пользу следующего тезиса:

\begin{quote}
узел может быть широким в двух разных смыслах: потому что у него много исходящих
каналов и потому что он сам по себе представляет более объемное понятие.
\end{quote}

Именно эта двойственность и пыталась поймать пара ``capacity + leak''.

\begin{riskbox}
\texttt{exp6b} остается синтетическим опытом. Capacity в нем ручная. Поэтому
после него нельзя было честно говорить ``мы нашли закон памяти''. Можно было
говорить только ``семейство без capacity уже не выглядит достаточным''.
\end{riskbox}

\section{\texttt{exp6c}: mechanism benchmark для знаменателей}

\subsection{Зачем понадобился еще один шаг}

У \texttt{exp6b} оставалось методологическое возражение: возможно, ответ
приходил не столько из-за хорошего знаменателя, сколько из-за удачной geometry
anchors и semantic gate. Поэтому \texttt{exp6c} радикально очищает постановку:

\begin{itemize}[leftmargin=2em]
\item для каждого запроса заранее задан source event;
\item после этого диффузия идет без query-conditioned gating;
\item benchmark измеряет именно поведение знаменателя.
\end{itemize}

\subsection{Диагностические мотивы}

Здесь уже не большой смешанный мир, а набор маленьких мотивов, каждый из которых
стрессует отдельное различие:

\begin{itemize}[leftmargin=2em]
\item weather motif -- проверка weighted leakage;
\item kit motif -- наказание слишком грубого degree counting;
\item travel motif -- чистый capacity-only case;
\item health motif -- перенос исходного \texttt{caffeine}/\texttt{beverage}
конфликта в новый стенд.
\end{itemize}

\subsection{Результат}

В \texttt{exp6c} baseline и degree-control выступили слабо:

\begin{itemize}[leftmargin=2em]
\item $R_0$ решил 2/5 задач;
\item $R_1$ в leave-one-out вообще рухнул до 1/5;
\item weighted leak alone ($R_2$, $R_3$) улучшил ситуацию, но не дотянул до
полной стабильности;
\item capacity-aware families ($R_4$, $R_5$) оказались единственными по-настоящему
сильными кандидатами.
\end{itemize}

Лучшим leave-one-out вариантом стал

\[
R_5(B)=1.0+\capacity(B)+0.05\, \sout(B)^{0.5}.
\]

\subsection{Что добавил \texttt{exp6c} относительно \texttt{exp6b}}

\texttt{exp6b} говорил: ``capacity выглядит полезной''. \texttt{exp6c} добавил:
``capacity делает свою работу даже тогда, когда anchor geometry убрана из игры''.

Это уже не просто победа в смешанном benchmark, а настоящий mechanism audit.

\section{\texttt{exp6d}: глобальные priors поверх локального удержания}

\subsection{Почему после \texttt{6c} вопрос еще не был закрыт}

Даже если узел хорошо удерживает локальный приток, он может быть ярким просто
потому, что глобально структурно знаменит. Поэтому возникает следующий шаг:

\[
\score(v\mid q)=\frac{\text{LocalResponse}(v\mid q)}{\text{Prior}(v)^\eta}.
\]

Содержательный вопрос уже не в том, ``заменяет ли PageRank локальную физику'',
а в том, добавляет ли глобальный prior \emph{нечто не сводимое} к локальному
удержанию.

\subsection{Методы}

Сравнивались:

\begin{itemize}[leftmargin=2em]
\item $P0$: raw diffusion;
\item $P1$: raw / out-strength$^\eta$;
\item $P2$: raw / in-strength$^\eta$;
\item $P3$: raw / PageRank$^\eta$;
\item $P4$: local-only retained diffusion;
\item $P5$: local retained diffusion / PageRank$^\eta$.
\end{itemize}

При этом локальная часть в $P4$ и $P5$ фиксировалась на победителе из
\texttt{6c}.

\subsection{Ключевой мотив: ложная глобальная слава}

Особенно важным был morning motif, где \texttt{water} делали глобально известным
за счет множества дальних входящих связей, но локальный утренний cue делали почти
симметричным между \texttt{water} и \texttt{espresso}. Этот мотив проверял,
может ли prior отличить:

\begin{itemize}[leftmargin=2em]
\item ``узел яркий, потому что глобально везде встречается'';
\item ``узел яркий, потому что именно здесь он нужен''.
\end{itemize}

\subsection{Результат}

Лучшим методом стал гибрид

\[
P5(v\mid q)=\frac{\text{LocalResponse}(v\mid q)}{\PR(v)^{0.25}},
\]

где LocalResponse задается capacity-aware local denominator из \texttt{6c}.

Этот гибрид решил все 7 задач и in-sample, и leave-one-out.

\subsection{Что особенно важно в этом результате}

Успех \texttt{P5} означает не ``PageRank победил'', а нечто более тонкое:

\begin{enumerate}[leftmargin=2em]
\item purely global penalties полезны;
\item purely local retention полезна;
\item но наилучшее поведение дает их комбинация;
\item при этом $\eta$ должен быть мягким: полная PMI-подобная сила оказывается
слишком агрессивной.
\end{enumerate}

\begin{takeaway}
К концу серии \texttt{6x} складывается первая действительно зрелая retrieval
картина:
\[
\score(v\mid q)\approx
\frac{\incoming(v\mid q)}{R_{\mathrm{local}}(v)\,\PR(v)^\eta}.
\]
Локальная часть отвечает за удержание и противодействие smearing, глобальная --
за снятие structural fame. Это еще не перенос на реальный граф, но это уже
чистая и достаточно сильная рабочая гипотеза.
\end{takeaway}

\section{Что серия \texttt{6x} дала и чего не дала}

\subsection{Что она действительно установила}

\begin{enumerate}[leftmargin=2em]
\item Пост-аккумуляционная topology-aware correction имеет смысл.
\item Простая степень узла может работать как грубый control, но она слишком
тупа, чтобы считаться итоговым законом.
\item Weighted leak полезен, но сам по себе не исчерпывает задачу.
\item Capacity выполняет не декоративную, а содержательную работу.
\item Глобальный prior, особенно мягкий PageRank, добавляет нечто реальное поверх
локальной физики.
\end{enumerate}

\subsection{Чего серия еще не установила}

\begin{enumerate}[leftmargin=2em]
\item Что manual capacity переносима на реальный граф.
\item Что найденная формула не является артефактом hand-authored motifs.
\item Что hybrid local + global сохранит преимущество на шумном эмбриональном
подграфе.
\end{enumerate}

Именно поэтому после \texttt{6d} естественным следующим ходом становится не
очередной synthetic benchmark, а transfer-check на реальном seeded subgraph.
